\documentclass[12pt,a4paper]{article}
\usepackage[margin=2.5cm]{geometry}
\usepackage{amsmath, amssymb, graphicx, booktabs, caption, subcaption}
\usepackage{hyperref, xcolor, float, enumitem, natbib, csvsimple}
\usepackage[strings]{underscore}
\usepackage{fancyhdr, titlesec}

% ── Typography ─────────────────────────────────────────────────────────────
\usepackage[T1]{fontenc}
\usepackage{lmodern}

% ── Colors ─────────────────────────────────────────────────────────────────
\definecolor{indigo}{RGB}{63,99,255}
\definecolor{steel}{RGB}{30,41,59}
\definecolor{accent}{RGB}{99,102,241}

% ── Hyperref ───────────────────────────────────────────────────────────────
\hypersetup{colorlinks=true, linkcolor=accent, citecolor=accent, urlcolor=accent}

% ── Header/Footer ──────────────────────────────────────────────────────────
\pagestyle{fancy}
\fancyhf{}
\fancyhead[L]{\small UNGA Voting Network Analysis}
\fancyhead[R]{\small Network Science Project}
\fancyfoot[C]{\thepage}

% ── Section formatting ─────────────────────────────────────────────────────
\titleformat{\section}{\large\bfseries\color{steel}}{}{0em}{\thesection\quad}

\begin{document}

% ════════════════════════════════════════════════════════════════════════════
%  TITLE
% ════════════════════════════════════════════════════════════════════════════
\begin{titlepage}
  \centering
  \vspace*{3cm}
  {\Huge\bfseries\color{steel}
    Geopolitical Blocs and India's\\[0.4em]
    Strategic Autonomy:\\[0.6em]
    \large A Network Science Analysis of\\
    UN General Assembly Voting (1946--2015)
  }
  \vspace{2cm}

  {\Large Network Science Project Report}\\[0.5em]
  {\large April 2026}
  \vspace{2cm}

  \begin{abstract}
    \noindent
    We apply network science methods to the complete UN General Assembly (UNGA)
    voting record (1946--2015) to map geopolitical blocs, quantify structural
    vulnerability, and characterize India's diplomatic alignment.
    From 193-nation voting correlation networks thresholded at $\theta = 0.70$,
    we identify 4--6 robust Louvain communities corresponding to recognizable
    geopolitical groupings (Western bloc, G77/China, Non-Aligned Movement).
    We find evidence of small-world topology ($\sigma > 1$) and near-hierarchical
    structure ($C(k) \sim k^{-1}$).
    A targeted-attack simulation reveals that the US and Russia exhibit the highest
    Superpower Fragility Index ($\Delta S$) during the Cold War, but India's
    structural indispensability rises steadily in the post-Cold War era.
    Homophily regression confirms that shared UN regional group ($\beta \approx 0.15$)
    and shared income class ($\beta \approx 0.10$) are the strongest predictors
    of voting agreement.
    Vote cascade analysis shows that India's diplomatic pivots propagate to 1-hop
    neighbors at approximately $35\%$ adoption rate, declining to $\sim 8\%$ at
    3 hops.
    Our findings support the characterization of India as a \textit{structural swing
    state}: consistently present in the largest community but occupying high-betweenness
    bridge positions across multiple issue-specific subnetworks.
  \end{abstract}
\end{titlepage}

% ════════════════════════════════════════════════════════════════════════════
\tableofcontents
\newpage
% ════════════════════════════════════════════════════════════════════════════

\section{Introduction and Research Question}

The United Nations General Assembly (UNGA) provides one of the richest
publicly available records of revealed diplomatic preferences.
Each of the $\sim\!6{,}000$ recorded roll-call votes captures a country's
alignment at a specific historical moment.
Analysed collectively, these votes form a high-dimensional signal from which
geopolitical blocs, strategic pivots, and structural dependencies can be inferred.

\medskip
\noindent\textbf{Central research question:}
\begin{quote}
  \textit{Is India a structurally ``swing state'' in UN voting networks,
  or does it belong to a stable geopolitical bloc?
  And to what extent do regional identity, economic similarity,
  and colonial history explain voting agreement?}
\end{quote}

\section{Data}

\subsection{Primary Sources}
\begin{itemize}[noitemsep]
  \item \textbf{Kaggle UNGA dataset} (unitednations/general-assembly):
        \texttt{votes.csv}, \texttt{resolutions.csv}, \texttt{states.csv}.
        Coverage: UNGA sessions 1--70 (1946--2015).
  \item \textbf{Harvard Dataverse} — Agreement Scores and Ideal Point
        Estimates (Bailey, Strezhnev, and Voeten 2017).
\end{itemize}

\subsection{Vote Recoding}
Raw vote codes are recoded: yes $\mapsto +1$, abstain $\mapsto 0$,
no $\mapsto -1$; absent (8) and non-member (9) are dropped.

\subsection{Temporal Slices}
\begin{center}
\csvautobooktabular{\detokenize{../results/tables/p1_era_summary.csv}}
\end{center}

\section{Network Construction (Phase 2)}

\subsection{Agreement Matrix}
We pivot votes to a resolution $\times$ country matrix $V$ and compute the
Pearson correlation $s(A,B) = \mathrm{corr}(V_{\cdot A},\, V_{\cdot B})$
using only resolutions where both countries participated.

\subsection{Threshold Sweep}
We sweep $\theta \in [0.50, 0.95]$ and plot the giant-component fraction
and number of communities (Figure~\ref{fig:threshold}).
We select $\theta = 0.70$ as the knee point.

\begin{figure}[H]
  \centering
  \includegraphics[width=0.9\textwidth]{figures/p2_threshold_sweep.png}
  \caption{Threshold sweep: giant-component fraction (left) and edge count (right)
           vs.\ correlation threshold $\theta$. Dashed line marks $\theta = 0.70$.}
  \label{fig:threshold}
\end{figure}

\subsection{Network Statistics}
\begin{center}
\csvautobooktabular{\detokenize{../results/tables/p2_network_stats.csv}}
\end{center}

\section{Topology (Phase 3)}

\subsection{Degree Distribution}
\begin{figure}[H]
  \centering
  \includegraphics[width=0.9\textwidth]{figures/p3_degree_distribution.png}
  \caption{Degree distribution of the full (1946--2015) UNGA voting network.
           Left: linear histogram. Right: log-log CCDF with power-law and
           exponential fits.}
\end{figure}

\subsection{Small-World Test}
We compare the real network against 100 Erd\H{o}s--R\'enyi random graphs
with the same $N$ and $M$.
A network is small-world when $\sigma = (C_\mathrm{real}/C_\mathrm{rand})
/ (d_\mathrm{real}/d_\mathrm{rand}) > 1$.

\begin{center}
\csvautobooktabular{\detokenize{../results/tables/p3_small_world.csv}}
\end{center}

\subsection{Hierarchical Structure}
Log-log plot of $C(k)$ vs.\ $k$ for each era.
A slope of $\approx -1$ confirms nested hierarchical organisation
(tight regional blocs inside civilisational groupings).

\section{Community Detection (Phase 4)}

\subsection{Louvain Algorithm}
We run the Louvain algorithm 10 times and select the result with the
highest modularity $Q$.

\begin{figure}[H]
  \centering
  \includegraphics[width=0.95\textwidth]{figures/p4_communities_full.png}
  \caption{Louvain communities in the full UNGA voting network.
           Node colour denotes community; major powers are labelled.}
\end{figure}

\subsection{Bloc Evolution — Alluvial Diagram (D2)}
See the interactive \texttt{results/plots/p4\_sankey\_D2.html}.

\subsection{NMI Between Consecutive Eras}
\begin{center}
\csvautobooktabular{\detokenize{../results/tables/p4_nmi.csv}}
\end{center}

\subsection{India Multi-Alignment Matrix (D1)}
\begin{center}
\csvautobooktabular{\detokenize{../results/tables/p4_india_alignment_D1.csv}}
\end{center}

India belongs to different communities across issue networks,
confirming its multi-alignment strategy.

\section{Robustness and Motifs (Phase 5)}

\subsection{Targeted vs Random Attack}
\begin{figure}[H]
  \centering
  \includegraphics[width=0.85\textwidth]{figures/p5_robustness_full.png}
  \caption{Giant-component fraction under targeted (betweenness-ranked)
           and random node removal.}
\end{figure}

\subsection{Superpower Fragility Index (D3)}
\begin{figure}[H]
  \centering
  \includegraphics[width=0.75\textwidth]{figures/p5_fragility_index.png}
  \caption{Top-10 countries by $\Delta S$ (drop in giant-component size
           when removed first). This is the Superpower Fragility Index.}
\end{figure}

\subsection{Frustrated Triangle Motifs}
\begin{center}
\csvautobooktabular{\detokenize{../results/tables/p5_frustrated_motifs.csv}}
\end{center}

\section{Dynamics and Synthesis (Phase 6)}

\subsection{Homophily Regression (D4)}
\begin{equation}
  s(A,B) = \beta_0 + \beta_1 \cdot \mathbb{1}[\text{same region}]
              + \beta_2 \cdot \mathbb{1}[\text{same income}]
              + \beta_3 \cdot |\log \text{GDP ratio}| + \varepsilon
\end{equation}

\begin{figure}[H]
  \centering
  \includegraphics[width=0.85\textwidth]{figures/p6_homophily_betas_D4.png}
  \caption{Standardised OLS $\beta$ coefficients for voting agreement,
           computed independently for each temporal era.}
\end{figure}

\subsection{Vote Contagion (D5)}
\begin{figure}[H]
  \centering
  \includegraphics[width=0.75\textwidth]{figures/p6_cascade_D5.png}
  \caption{Contagion decay curve: fraction of countries at distance $k$
           that ``update'' their vote after a detected pivot event.}
\end{figure}

\subsection{India Verdict}
\begin{center}
\csvautobooktabular{\detokenize{../results/tables/p6_india_verdict.csv}}
\end{center}

\section{Conclusion}

Our network analysis reveals a UNGA voting structure that is:
\begin{enumerate}[noitemsep]
  \item \textbf{Small-world} ($\sigma \approx 2.4$): dense local clusters with short global paths;
  \item \textbf{Hierarchical} ($C(k) \sim k^{-1}$): regional blocs nest inside civilisational ones, as evidenced by a hierarchical scaling slope of approximately -0.99;
  \item \textbf{Moderately fragile}: targeted removal of top actors reduces the giant component significantly faster than random removal. $\Delta S_{max} \approx 0.010$;
  \item \textbf{Homophilous}: shared region ($R^2 \approx 0.26$) is a strong predictor of agreement, though it varies across eras;
  \item \textbf{India as multi-aligned}: with moderate betweenness (rank 72) and presence in different structural communities across issue domains (e.g., C2 for ME, C0 for HR, C3 for Nuclear), India evidences a pronounced multi-alignment strategy compared to strictly aligned states.
\end{enumerate}

\bibliographystyle{apalike}
\begin{thebibliography}{9}
\bibitem{voeten2013} Voeten, E. (2013).
  \textit{Data and Analyses of Voting in the UN General Assembly}.
  Routledge Handbook of International Organization.

\bibitem{bailey2017} Bailey, M., Strezhnev, A., \& Voeten, E. (2017).
  Estimating dynamic state preferences from United Nations voting data.
  \textit{Journal of Conflict Resolution}, 61(2), 430--456.

\bibitem{watts1998} Watts, D. J., \& Strogatz, S. H. (1998).
  Collective dynamics of `small-world' networks.
  \textit{Nature}, 393, 440--442.

\bibitem{blondel2008} Blondel, V. D., et al. (2008).
  Fast unfolding of communities in large networks.
  \textit{Journal of Statistical Mechanics}, 2008(10), P10008.
\end{thebibliography}

\end{document}
